% =====================================================================
%  Lecture 02 -- 2G Circuit Switching                     (Persian edition)
% =====================================================================
\documentclass[11pt,a4paper]{article}
\providecommand\ACNroot{../..}
\input{\ACNroot/templates/lecture-fa}

\ACNset{number}{02}
\ACNset{title}{سوئیچ مداری در نسل دوم}
\ACNset{shorttitle}{سوئیچ مداری در نسل دوم}
\ACNset{titleother}{2G Circuit Switching}
\ACNset{part}{بخش اول -- مبانی نسل دوم و سوم (درس‌های ۱ تا ۴)}
\ACNset{time}{۶ تا ۸ ساعت (۳ ساعت مطالعه، ۳ تا ۵ ساعت تمرین)}
\ACNset{level}{متوسط}
\ACNset{prereq}{درس ۱ (معماری \lr{GSM})؛ احتمال پایه (فرایند پواسون، توزیع نمایی)}
\ACNset{gen}{2}

\begin{document}
\ACNcover

% ---------------------------------------------------------------------
\section{اهداف یادگیری}
پس از مطالعهٔ این درس باید بتوانید:
\begin{objectives}
  \item \textbf{توضیح دهید} مدار چیست، یک تماس \acro{GSM} چه منابعی را رزرو می‌کند و چرا
        رزرو کیفیت را تضمین می‌کند اما ظرفیت را هدر می‌دهد.
  \item پشتهٔ پروتکل \acro{SS7} (\lr{MTP, SCCP, TCAP, MAP, ISUP, BSSAP}) و نقش نقاط سیگنالینگ و
        نقاط انتقال سیگنال را \textbf{توصیف کنید}.
  \item یک تماس آغازشده از موبایل و یک تماس پایان‌یافته به موبایل را، شامل پرسش از \acro{HLR}
        و استفاده از \acro{MSRN}، \textbf{دنبال کنید}.
  \item فرمول \lr{Erlang-B} را از مدل $M/M/m/m$ \textbf{استخراج کنید} و یک سلول یا گروه ترانک
        را برای احتمال انسداد هدف \textbf{ابعادبندی کنید}.
  \item بازده ترانکینگ و هزینهٔ تقسیم یک مخزن کانال را \textbf{کمّی کنید}.
  \item سیگنالینگ رومینگ و مسیریابی تماس را \textbf{تحلیل کنید} و \textbf{ارزیابی کنید} چرا
        سوئیچ مداری برای داده به یک محدودیت تبدیل شد.
\end{objectives}

% ---------------------------------------------------------------------
\section{دانش پیش‌نیاز}
به معماری و هویت‌های \acro{GSM} از درس ۱ (\lr{MSC, VLR, HLR, GMSC, IMSI, TMSI}) و احتمال
مقدماتی نیاز دارید: فرایند ورود پواسون با نرخ $\lambda$ و زمان اشغال با توزیع نمایی و
میانگین $h$. واژهٔ \emph{ترافیک} را به معنای تلفنی آن به کار می‌بریم: میانگین تعداد مدارهای
هم‌زمان مشغول.

\paragraph{واژه‌های کلیدی} واژگان این درس همراه با معادل انگلیسی آن‌ها:
\begin{keyterms}
سوئیچ مداری & \lr{Circuit switching} & یک مسیر اختصاصی سرتاسری برای کل مدت تماس رزرو می‌شود. \\
سامانهٔ سیگنالینگ شمارهٔ ۷ & \lr{SS7} & شبکهٔ سیگنالینگ خارج از باند میان مراکز سوئیچ. \\
بخش کاربر \acro{ISDN} & \lr{ISUP} & پروتکل \acro{SS7} برای برقراری و آزادسازی مدارهای ترانک. \\
بخش کاربرد موبایل & \lr{MAP} & پروتکل کاربردی \acro{SS7} برای تحرک (\lr{HLR, VLR, MSC}). \\
شمارهٔ تلفن مشترک & \lr{MSISDN} & شماره‌ای که برای تماس با مشترک گرفته می‌شود. \\
شمارهٔ رومینگ ایستگاه موبایل & \lr{MSRN} & شمارهٔ موقتی برای مسیریابی تماس به \acro{MSC} خدمت‌دهنده. \\
ارلانگ & \lr{Erlang (E)} & واحد ترافیک: یک مدار که پیوسته مشغول است. \\
احتمال انسداد & \lr{Blocking probability} & احتمال اینکه یک تلاش تماس همهٔ مدارها را مشغول بیابد. \\
درجهٔ خدمت & \lr{Grade of service} & احتمال انسداد هدف که در ابعادبندی به کار می‌رود. \\
بازده ترانکینگ & \lr{Trunking efficiency} & ترافیک حمل‌شده به ازای هر مدار؛ با اندازهٔ مخزن رشد می‌کند. \\
رومینگ & \lr{Roaming} & استفاده از شبکهٔ میزبان با اشتراک خانگی. \\
\end{keyterms}

% ---------------------------------------------------------------------
\section{مقدمه و انگیزه}
تلفن بر یک وعدهٔ ساده بنا شد: وقتی طرف مقابل پاسخ داد، مسیر صدا مانند یک جفت سیم رفتار
می‌کند، با تأخیر ثابت، بدون اتلاف و بدون رقابت با کاربران دیگر. راه مهندسی برای عمل به این
وعده \emph{رزرو} منابع برای کل مدت تماس است. \acro{GSM} این مدل را از \acro{ISDN} به ارث برد،
زیرا نخستین وظیفه‌اش گسترش موبایل شبکهٔ تلفن ثابت بود و مراکز سوئیچ، ترانک‌ها و سامانه‌های
صورت‌حساب دههٔ ۱۹۹۰ همگی بر پایهٔ مدارهای \qty{64}{\kilo\bit\per\second} ساخته شده بودند.

دو مسئله باید افزون بر شبکهٔ تلفن کلاسیک حل می‌شد. نخست، \emph{سیگنالینگ} باید چیزی بسیار
بیش از ارقام حمل می‌کرد: احراز هویت، به‌روزرسانی مکان، تحویل و رومینگ میان کشورها. دوم،
شبکه باید پیش از برقراری هر مداری مشترکی را که مکانش تغییر می‌کند \emph{پیدا} می‌کرد.
\acro{GSM} هر دو را با سامانهٔ سیگنالینگ شمارهٔ ۷ (\acro{SS7}) و بخش کاربرد موبایل آن
(\acro{MAP}) حل کرد. این درس یک تماس را از شماره‌گیری تا قطع دنبال می‌کند و سپس پرسش کمّی‌ای
را مطرح می‌کند که هر اپراتوری با آن روبه‌روست: \emph{چند مدار کافی است؟}

% ---------------------------------------------------------------------
\section{مفاهیم بنیادی}

\subsection{مدار چه چیزی را رزرو می‌کند}
برای یک تماس صوتی \acro{GSM}، شبکه در کل مدت تماس این منابع را رزرو می‌کند:
\begin{itemize}
  \item یک کانال ترافیک (\acro{TCH}) روی واسط رادیویی (یک شیار زمانی روی یک حامل)؛
  \item یک زیرکانال \qty{16}{\kilo\bit\per\second} روی پیوند \acro{Abis} و یک شیار
        \qty{64}{\kilo\bit\per\second} روی واسط \acro{A} پس از تبدیل کدک؛
  \item یک اتصال در ماتریس سوئیچ \acro{MSC}؛
  \item یک ترانک \qty{64}{\kilo\bit\per\second} روی هر پیوند میان مراکز (یک مدار \acro{ISUP} در هر پیوند).
\end{itemize}
این رزرو حتی در سکوت حفظ می‌شود. در یک مکالمهٔ معمولی هر جهت تنها حدود ۴۰٪ زمان فعال است؛
بنابراین بیشتر ظرفیت رزروشده چیزی حمل نمی‌کند.

\begin{definitionbox}{شدت ترافیک}
اگر تماس‌ها با نرخ $\lambda$ (تماس در ثانیه) برسند و هر یک به‌طور میانگین $h$ ثانیه طول بکشد،
ترافیک ارائه‌شده $A = \lambda h$ است؛ عددی بی‌بعد که بر حسب ارلانگ اندازه‌گیری می‌شود. یک ارلانگ
ترافیکی است که یک مدار را همیشه مشغول نگه دارد.
\end{definitionbox}

\subsection{سیگنالینگ درون‌باند در برابر سیگنالینگ کانال مشترک}
شبکه‌های تلفن اولیه با تُن‌هایی روی همان مدار صوتی سیگنالینگ می‌کردند. این روش کند بود، واژگان
محدودی داشت و ناامن بود (کاربران می‌توانستند تُن‌ها را تقلید کنند). سیگنالینگ کانال مشترک این
دو را جدا می‌کند: پیام‌های سیگنالینگ روی یک شبکهٔ بسته‌ای اختصاصی حرکت می‌کنند و مدارهای صوتی
تنها صدا حمل می‌کنند. \acro{SS7} استاندارد \lr{ITU-T} برای این شبکهٔ سیگنالینگ است \cite{ituq700}.
این نخستین نمونه در این درس از \emph{جدایی صفحهٔ کنترل و صفحهٔ کاربر} است؛ ایده‌ای که در هر
نسل بازمی‌گردد.

\subsection{نقاط سیگنالینگ}
گره‌های شبکهٔ \acro{SS7} \emph{نقاط سیگنالینگ} (\acro{SP}) هستند که هر یک با یک کد نقطه شناسایی
می‌شود. \emph{نقطهٔ سوئیچ خدمت} (\acro{SSP}) یک مرکز سوئیچ مانند \acro{MSC} است. \emph{نقطهٔ
انتقال سیگنال} (\acro{STP}) مسیریاب بسته‌ای شبکهٔ \acro{SS7} است: پیام‌ها را میان \acro{SP}ها
بدون پایان‌دادن به آن‌ها جابه‌جا می‌کند و برای قابلیت اطمینان به‌صورت جفت‌های همزاد نصب
می‌شود. \emph{نقطهٔ کنترل خدمت} (\acro{SCP}) پایگاه‌های داده و منطق خدمت را در بر دارد، برای
مثال منطق شبکهٔ هوشمند برای صورت‌حساب پیش‌پرداخت.

% ---------------------------------------------------------------------
\section{توضیحات فنی تفصیلی}

\subsection{پشتهٔ پروتکل \acro{SS7}}
\Cref{fig:ss7} پشته را آن‌گونه که در \acro{GSM} به کار می‌رود نشان می‌دهد.
\begin{itemize}
  \item \textbf{\lr{MTP 1--3}} (بخش انتقال پیام) پیام‌ها را به‌طور قابل اطمینان میان گره‌های
        مجاور حمل می‌کند (\lr{MTP 2}: قاب‌بندی، آشکارسازی خطا و ارسال مجدد) و آن‌ها را با کد
        نقطه در شبکهٔ سیگنالینگ مسیریابی می‌کند (\lr{MTP 3}).
  \item \textbf{\acro{SCCP}} (بخش کنترل اتصال سیگنالینگ) خدمات بی‌اتصال و اتصال‌گرا و، مهم‌تر
        از همه برای تحرک، \emph{ترجمهٔ عنوان سراسری} را می‌افزاید: پیامی را می‌توان با یک عنوان
        سراسری شبیه شمارهٔ تلفن (برای مثال برگرفته از \acro{IMSI}) نشانی‌دهی کرد و شبکه کد نقطهٔ
        \acro{HLR} را پیدا می‌کند.
  \item \textbf{\acro{TCAP}} (بخش کاربرد قابلیت‌های تراکنش) گفت‌وگوهای درخواست--پاسخ میان
        کاربردها را ساختار می‌دهد.
  \item \textbf{\acro{MAP}} \cite{ts29002} روی \acro{TCAP} اجرا می‌شود و عملیات تحرک را تعریف
        می‌کند: \msg{Update Location}، \msg{Insert Subscriber Data}، \msg{Send Authentication Info}،
        \msg{Send Routing Information} (\acro{SRI})، \msg{Provide Roaming Number} (\acro{PRN}) و
        \msg{Cancel Location}.
  \item \textbf{\acro{ISUP}} (بخش کاربر \acro{ISDN}) مستقیماً روی \lr{MTP 3} اجرا می‌شود و مدارهای
        ترانک میان مراکز را با پیام‌هایی مانند \lr{IAM, ACM, ANM, REL, RLC} برقرار و آزاد می‌کند.
  \item \textbf{\acro{BSSAP}} روی \acro{SCCP} و در واسط \acro{A} اجرا می‌شود. بخش \acro{BSSMAP} آن
        پیام‌های میان \acro{MSC} و \acro{BSC} را حمل می‌کند (تخصیص، تحویل، رمزنگاری)؛ بخش
        \acro{DTAP} آن پیام‌های میان \acro{MSC} و \acro{MS} را به‌طور شفاف حمل می‌کند (کنترل تماس
        و مدیریت تحرک).
\end{itemize}

\begin{acnfigure}
  \acnfig{lecture-02/ss7-stack}
  \caption{پشتهٔ \acro{SS7} در شبکهٔ هستهٔ \acro{GSM}. \acro{MTP} انتقال را فراهم می‌کند؛ \acro{ISUP}
  مدارهای ترانک را کنترل می‌کند؛ \acro{SCCP} و \acro{TCAP} از گفت‌وگوهای تحرک \acro{MAP} و
  \acro{BSSAP} در واسط \acro{A} پشتیبانی می‌کنند.}
  \label{fig:ss7}
\end{acnfigure}

\subsection{پروتکل‌های لایهٔ ۳ میان \acro{MS} و \acro{MSC}}
در سمت رادیو، لایهٔ ۳ به سه زیرلایه تقسیم می‌شود \cite{ts24008}: \emph{مدیریت منابع رادیویی}
(\acro{RR}، که در \acro{BSS} پایان می‌یابد: تخصیص کانال، تحویل)، \emph{مدیریت تحرک}
(\acro{MM}: به‌روزرسانی مکان، احراز هویت، تخصیص مجدد \acro{TMSI}) و \emph{مدیریت اتصال}
(\acro{CM}، شامل کنترل تماس \acro{CC}: \lr{Setup, Alerting, Connect, Disconnect}). پیام‌های
\acro{MM} و \acro{CC} به‌طور شفاف در \acro{DTAP} از \acro{BSS} عبور می‌کنند.

\subsection{شماره‌ها و هویت‌های به‌کاررفته در مسیریابی}
سه شماره را نباید با هم اشتباه گرفت:
\begin{itemize}
  \item \textbf{\acro{MSISDN}} -- شماره‌ای که مردم می‌گیرند؛ به شبکهٔ خانگی مشترک اشاره
        می‌کند و \acro{GMSC} از آن برای یافتن \acro{HLR} استفاده می‌کند.
  \item \textbf{\acro{IMSI}} -- هویت اشتراک که درون شبکهٔ موبایل به کار می‌رود (نشانی‌دهی
        \acro{MAP}، احراز هویت)؛ هرگز شماره‌گیری نمی‌شود.
  \item \textbf{\acro{MSRN}} -- شمارهٔ موقت \lr{E.164} که \acro{VLR} میزبان برای یک تماس ورودی
        تخصیص می‌دهد؛ تماس را از طریق شبکهٔ ثابت به \acro{MSC} خدمت‌دهنده مسیریابی می‌کند.
\end{itemize}
\acro{HLR} شمارهٔ \acro{MSISDN} را به \acro{IMSI} نگاشت می‌کند و ذخیره می‌کند کدام \acro{VLR} به
مشترک خدمت می‌دهد؛ \acro{VLR} هنگام رسیدن تماس، \acro{MSRN} را دوباره به \acro{IMSI} نگاشت می‌کند.

\subsection{به‌روزرسانی مکان}
سه حالت به‌روزرسانی مکان را آغاز می‌کنند \cite{ts23012}:
\begin{itemize}
  \item \textbf{عادی} -- \acro{MS} شناسهٔ ناحیهٔ مکانی (\acro{LAI}) جدیدی روی کانال پخش می‌خواند.
        اگر \acro{LA} جدید متعلق به \acro{VLR} دیگری باشد، \acro{VLR} جدید داده‌های مشترک را با
        \msg{Update Location} از \acro{HLR} دریافت می‌کند و \acro{HLR} پیام \msg{Cancel Location} را
        به \acro{VLR} قبلی می‌فرستد.
  \item \textbf{دوره‌ای} -- \acro{MS} با پایان زمان‌سنج \lr{T3212} دوباره ثبت‌نام می‌کند تا شبکه
        بتواند مشترکانی را که بدون جداشدن ناپدید شده‌اند (باتری برداشته‌شده، از دست رفتن پوشش)
        پاک کند.
  \item \textbf{اتصال \acro{IMSI}} -- هنگام روشن‌شدن در همان \acro{LA}. رویهٔ قرینهٔ آن،
        \textbf{جداشدن \acro{IMSI}}، به \acro{VLR} اطلاع می‌دهد که \acro{MS} خاموش است تا تماس‌های
        ورودی بیهوده فراخوانی نشوند.
\end{itemize}

\subsection{رومینگ}
مشترک رومینگ در یک شبکهٔ میزبان (\acro{VPLMN}) ثبت‌نام می‌کند. \acro{VLR} میزبان با استفاده از
عنوان‌های سراسری برگرفته از \acro{IMSI}، از طریق شبکهٔ بین‌المللی \acro{SS7} به \acro{HLR} خانگی
(\acro{HPLMN}) دسترسی پیدا می‌کند و داده‌های اشتراک و سه‌تایی‌های احراز هویت را دریافت می‌کند.
توافق‌نامه‌های رومینگ صورت‌حساب و خدمات مجاز را تعیین می‌کنند. دو پیامد شایان توجه است:
\begin{itemize}
  \item \textbf{احراز هویت در دست شبکهٔ خانگی می‌ماند}: تنها \acro{AuC} خانگی $K_i$ را می‌داند.
  \item \textbf{مسیر رفت‌وبرگشتی (\lr{tromboning})}: تماس با یک مشترک رومینگ همیشه ابتدا به
        \acro{GMSC} خانگی (شمارهٔ \acro{MSISDN} متعلق به شبکهٔ خانگی است) و سپس به‌صورت
        بین‌المللی به \acro{MSC} میزبان مسیریابی می‌شود. بنابراین تماس‌گیرندهٔ محلی در کشور
        میزبان هزینهٔ دو مسیر بین‌المللی را می‌پردازد، مگر اینکه مسیریابی بهینه پیاده شده باشد.
\end{itemize}

\subsection{آزادسازی تماس}
هر یک از دو طرف می‌تواند تماس را آزاد کند. موبایل \msg{Disconnect} می‌فرستد؛ \acro{MSC} با
\msg{Release} پاسخ می‌دهد و \acro{MS} با \msg{Release Complete} تأیید می‌کند. به سمت شبکهٔ ثابت،
\acro{MSC} پیام \acro{ISUP} \msg{REL} را می‌فرستد و \msg{RLC} را دریافت می‌کند که ترانک را آزاد
می‌کند. در پایان \acro{MSC} به \acro{BSS} دستور آزادسازی کانال رادیویی را می‌دهد
(\acro{BSSMAP} \msg{Clear Command}). رکوردهای صورت‌حساب با آزاد شدن مدار بسته می‌شوند؛ به همین
دلیل صورت‌حساب سوئیچ مداری به‌طور طبیعی بر پایهٔ زمان است.

% ---------------------------------------------------------------------
\section{مدل‌های ریاضی}

\subsection{سامانهٔ اتلاف و فرمول \lr{Erlang-B}}\label{sec:erlangb}
\paragraph{مسئلهٔ مهندسی} یک سلول $m$ کانال ترافیک دارد. تماس‌هایی که وقتی همه مشغول‌اند
می‌رسند \emph{مسدود} (از دست رفته) می‌شوند؛ منتظر نمی‌مانند. $m$ چقدر باید بزرگ باشد تا
بیش از، مثلاً، ۲٪ تلاش‌ها مسدود نشوند؟

\paragraph{مدل سامانه و فرض‌ها} (۱) ورود تماس‌ها یک فرایند پواسون با نرخ $\lambda$ از یک جمعیت
بزرگ است؛ (۲) زمان‌های اشغال مستقل و با توزیع نمایی و میانگین $h = 1/\mu$ هستند؛ (۳) تماس‌های
مسدودشده حذف می‌شوند (بدون تلاش مجدد)؛ (۴) همهٔ $m$ کانال به‌طور کامل در دسترس‌اند. این همان
صف $M/M/m/m$ است.

\paragraph{استخراج} فرض کنید $p_k$ احتمال حالت پایدار مشغول بودن $k$ کانال باشد. در حالت $k$
سامانه با نرخ $\lambda$ بالا و با نرخ $k\mu$ پایین می‌رود. برقراری توازن جریان میان حالت‌های
مجاور نتیجه می‌دهد
\begin{equation}\label{eq:balance}
  \lambda\,p_{k-1} = k\mu\,p_k \quad\Longrightarrow\quad
  p_k = \frac{A}{k}\,p_{k-1} = \frac{A^k}{k!}\,p_0, \qquad A = \frac{\lambda}{\mu} = \lambda h .
\end{equation}
از بهنجارش $\sum_{k=0}^{m}p_k = 1$ داریم $p_0 = \bigl(\sum_{i=0}^{m}A^i/i!\bigr)^{-1}$. چون
ورودهای پواسون میانگین‌های زمانی را می‌بینند (ویژگی \acro{PASTA})، احتمال مسدود شدن یک تماس
ورودی برابر احتمال مشغول بودن همهٔ $m$ کانال است:
\begin{equation}\label{eq:erlangb}
  P_{\mathrm{B}} = B(A,m) = \frac{A^m/m!}{\displaystyle\sum_{i=0}^{m}\frac{A^i}{i!}} .
\end{equation}
\paragraph{نمادها و یکاها} $A$: ترافیک ارائه‌شده [\lr{E}]؛ $m$: تعداد کانال‌ها؛ $\lambda$: نرخ
ورود [تماس بر ثانیه]؛ $h$: میانگین زمان اشغال [\lr{s}]؛ $P_{\mathrm{B}}$: احتمال انسداد (بی‌بعد).
ترافیک حمل‌شده $A_{\mathrm{c}} = A(1-P_{\mathrm{B}})$ و بهره‌وری هر کانال $\rho = A_{\mathrm{c}}/m$ است.

\paragraph{صورت پایدار عددی} محاسبهٔ مستقیم \cref{eq:erlangb} برای $m$ بزرگ سرریز می‌کند. با
تقسیم صورت و مخرج بر $A^m/m!$ رابطهٔ بازگشتی زیر به دست می‌آید
\begin{equation}\label{eq:erlangrec}
  B(A,0) = 1, \qquad B(A,k) = \frac{A\,B(A,k-1)}{k + A\,B(A,k-1)}, \quad k = 1,\dots,m,
\end{equation}
که در \lr{\texttt{scripts/acnmath.py}} برای تولید \cref{fig:erlang} و همهٔ اعداد زیر به کار رفته است.

\paragraph{تفسیر} \Cref{eq:erlangb} تنها به \emph{میانگین} زمان اشغال وابسته است (به توزیع زمان
اشغال حساس نیست) و به همین دلیل یک قرن برای تلفن به‌خوبی کار کرد. با این حال فرض می‌کند
کاربران مسدودشده می‌روند؛ در واقعیت بسیاری دوباره تلاش می‌کنند که ترافیک ارائه‌شده و انسداد
واقعی را افزایش می‌دهد (\cref{exr:retry}).

\begin{acnfigure}
  \acnfig{lecture-02/erlang-b}
  \caption{احتمال انسداد \lr{Erlang-B} بر حسب ترافیک ارائه‌شده برای $m = 8$، ۳۰ و ۹۴ کانال. برای
  درجهٔ خدمت یکسان ۲٪، هشت کانال \lr{3.6\,E} (\lr{0.45\,E} برای هر کانال) و ۹۴ کانال \lr{82.2\,E}
  (\lr{0.87\,E} برای هر کانال) حمل می‌کنند.}
  \label{fig:erlang}
\end{acnfigure}

\subsection{بازده ترانکینگ}
\paragraph{مسئله} آیا داشتن یک مخزن $2m$ کانالی بهتر است یا دو مخزن $m$ کانالی؟
\paragraph{نتیجه} برای درجهٔ خدمت ثابت، مخزن‌های بزرگ‌تر ترافیک بیشتری به ازای هر کانال حمل
می‌کنند. این \emph{بهرهٔ ترانکینگ} (یا بهرهٔ تجمیع) است و از \cref{eq:erlangb} نتیجه می‌شود:
نوسان نسبی تعداد تماس‌های فعال، $\sqrt{A}/A$، با رشد $A$ کوچک می‌شود و ظرفیت ذخیرهٔ کمتری
لازم است. \Cref{tab:trunking} آن را کمّی می‌کند. این توضیح می‌دهد چرا اپراتورهای \acro{GSM} از
تقسیم کانال‌ها میان بخش‌ها پرهیز می‌کنند مگر آنکه تداخل آن را لازم کند، و چرا تسهیم آماری در
شبکه‌های بسته‌ای این‌قدر جذاب است.

\begin{table}[htbp]
\centering\small
\caption{ترافیک حمل‌شده در انسداد ۲٪ (محاسبه‌شده با \cref{eq:erlangrec}).}\label{tab:trunking}
\begin{tabular}{@{}rrrr@{}}
\toprule
تعداد کانال $m$ & ترافیک ارائه‌شده $A$ [\lr{E}] & ترافیک حمل‌شده به ازای کانال [\lr{E}] & بهره‌وری \\
\midrule
8  & 3.63  & 0.44 & ۴۴٪ \\
15 & 9.01  & 0.59 & ۵۹٪ \\
30 & 21.93 & 0.72 & ۷۲٪ \\
47 & 37.46 & 0.78 & ۷۸٪ \\
94 & 82.17 & 0.86 & ۸۶٪ \\
\bottomrule
\end{tabular}
\end{table}

\subsection{هدررفت یک مدار رزروشده}
اگر ضریب فعالیت صدا $\alpha$ باشد (کسری از زمان که گوینده فعال است)، بار مفید یک مدار رزروشده
$\alpha$ برابر ظرفیت آن است. با $\alpha\approx0.4$ در هر جهت، حدود ۶۰٪ هر ثانیه-مدار رزروشده
سکوت حمل می‌کند. سوئیچ بسته‌ای این ظرفیت را به‌صورت آماری بازمی‌گرداند، به بهای تأخیر متغیر --
مصالحهٔ بنیادی میان دو الگوی سوئیچینگ که در درس‌های ۳ و ۵ دوباره بررسی می‌شود.

% ---------------------------------------------------------------------
\section{معماری شبکه و نمودارها}
شبکهٔ \acro{SS7} یک شبکهٔ بسته‌ای جداگانه است که روی شبکهٔ مداری قرار گرفته است. هر \acro{MSC}،
\acro{GMSC} و \acro{HLR} یک نقطهٔ سیگنالینگ است و \acro{STP}ها میان آن‌ها مسیریابی می‌کنند.
ترانک‌های میان \acro{MSC}ها و \acro{PSTN} صدا حمل می‌کنند، در حالی که پیام‌های \acro{MAP} و
\acro{ISUP} روی پیوندهای سیگنالینگی حرکت می‌کنند که ممکن است مسیرهای فیزیکی کاملاً متفاوتی
داشته باشند. \Cref{fig:mtc} این تعامل را برای سخت‌ترین حالت، یعنی تماس ورودی، و \cref{fig:moc}
حالت خروجی را نشان می‌دهد.

% ---------------------------------------------------------------------
\section{رویه‌های گام‌به‌گام}

\begin{procedurebox}{تماس آغازشده از موبایل (\acro{MOC})}\label{proc:moc}
\begin{steps}
  \item \acro{MS} درخواست کانال را روی کانال دسترسی تصادفی (\acro{RACH}) می‌فرستد و یک کانال
        سیگنالینگ (\acro{SDCCH}) را روی کانال اعطای دسترسی (\acro{AGCH}) دریافت می‌کند.
  \item پیام \msg{CM Service Request} را می‌فرستد؛ \acro{MSC/VLR} هویت \acro{MS} را احراز و
        رمزنگاری را آغاز می‌کند (درس ۱).
  \item \acro{MS} پیام \msg{Setup} را همراه با شمارهٔ مقصد می‌فرستد. \acro{MSC} اشتراک را بررسی
        می‌کند (محدودیت‌ها، اعتبار پیش‌پرداخت) و با \msg{Call Proceeding} پاسخ می‌دهد.
  \item \acro{MSC} به \acro{BSS} دستور تخصیص کانال ترافیک می‌دهد (\acro{BSSMAP}
        \msg{Assignment Request})؛ \acro{MS} از \acro{SDCCH} به یک \acro{TCH} می‌رود.
  \item \acro{MSC} شماره را تحلیل و \acro{ISUP} \msg{IAM} را به سوی مقصد ارسال می‌کند؛ مرکز
        دوردست هنگام زنگ خوردن تلفن مقصد \msg{ACM} را بازمی‌گرداند و \acro{MSC} پیام
        \msg{Alerting} را به \acro{MS} می‌فرستد (بوق انتظار).
  \item وقتی طرف مقابل پاسخ می‌دهد، \acro{ISUP} \msg{ANM} می‌رسد؛ \acro{MSC} پیام \msg{Connect}
        را می‌فرستد، \acro{MS} تأیید می‌کند و صورت‌حساب آغاز می‌شود.
\end{steps}
\end{procedurebox}

\begin{acnfigure}
  \resizebox{0.92\linewidth}{!}{\acnfig{lecture-02/moc-sequence}}
  \caption{تماس آغازشده از موبایل (\cref{proc:moc}).}
  \label{fig:moc}
\end{acnfigure}

\begin{procedurebox}{تماس پایان‌یافته به موبایل (\acro{MTC})}\label{proc:mtc}
\begin{steps}
  \item تماس‌گیرنده‌ای در \acro{PSTN} شمارهٔ \acro{MSISDN} را می‌گیرد. \acro{PSTN} تماس را
        (\acro{ISUP} \msg{IAM}) به یک \acro{GMSC} در شبکهٔ خانگی مسیریابی می‌کند.
  \item \acro{GMSC} پیام \acro{MAP} \msg{Send Routing Information} را به \acro{HLR} می‌فرستد.
  \item \acro{HLR} \acro{VLR} خدمت‌دهنده را پیدا می‌کند و با \acro{MAP} \msg{Provide Roaming Number}
        یک شمارهٔ رومینگ درخواست می‌کند؛ \acro{VLR} یک \acro{MSRN} تخصیص می‌دهد و بازمی‌گرداند.
  \item \acro{HLR} شمارهٔ \acro{MSRN} را به \acro{GMSC} برمی‌گرداند که تماس را با یک \msg{IAM}
        جدید به نشانی \acro{MSRN} مسیریابی می‌کند -- شبکهٔ ثابت آن را به \acro{MSC} خدمت‌دهنده
        تحویل می‌دهد.
  \item \acro{MSC} زمینهٔ مشترک را از \acro{VLR} خود می‌پرسد؛ \acro{VLR} شمارهٔ \acro{MSRN} را آزاد
        و ناحیهٔ مکانی را بازمی‌گرداند؛ \acro{MSC} موبایل را در همهٔ سلول‌های آن \acro{LA} فراخوانی می‌کند.
  \item \acro{MS} به فراخوانی پاسخ می‌دهد، احراز هویت می‌شود، \msg{Setup} و یک \acro{TCH} دریافت
        می‌کند؛ زنگ خوردن و پاسخ با \msg{ACM} و \msg{ANM} به عقب اعلام می‌شوند.
\end{steps}
\end{procedurebox}

\begin{acnfigure}
  \resizebox{\linewidth}{!}{\acnfig{lecture-02/mtc-sequence}}
  \caption{تماس پایان‌یافته به موبایل: پرسش از \acro{HLR}، تخصیص \acro{MSRN}، مسیریابی و فراخوانی
  (\cref{proc:mtc}). خط‌چین: سیگنالینگ؛ توپر: مدار رزروشده.}
  \label{fig:mtc}
\end{acnfigure}

% ---------------------------------------------------------------------
\section{مثال‌های حل‌شده}

\begin{examplebox}{ابعادبندی یک سلول}\label{ex:dim}
\textbf{مسئله.} سلول مثال ۱ درس ۱ دارای ۹۴ کانال ترافیک است. هر مشترک در ساعت اوج \lr{25\,mE}
ترافیک تولید می‌کند. این سلول با درجهٔ خدمت ۲٪ به چند مشترک خدمت می‌دهد؟

\textbf{حل.} حل $B(A,94) = 0.02$ با \cref{eq:erlangrec} مقدار $A = 82.2$\,E را می‌دهد. تعداد
مشترکان $N = A/a = 82.2/0.025 \approx 3290$ است. ترافیک حمل‌شده
$A_{\mathrm{c}} = 82.2\times0.98 = 80.5$\,E، یعنی بهره‌وری ۸۶٪ کانال‌ها.

\textbf{تفسیر.} اگر همین ۹۴ کانال میان دو بخش ۴۷تایی تقسیم می‌شد، هر بخش \lr{37.5\,E} و در مجموع
\lr{74.9\,E} حمل می‌کرد -- حدود ۹٪ مشترک کمتر با همان سخت‌افزار (\cref{tab:trunking}).
\end{examplebox}

\begin{examplebox}{محاسبهٔ دستی \lr{Erlang-B}}\label{ex:hand}
\textbf{مسئله.} یک سایت کوچک $m=3$ کانال و $A=2$\,E دارد. $P_{\mathrm{B}}$ را حساب کنید.

\textbf{حل.}
\begin{equation*}
  P_{\mathrm{B}} = \frac{2^3/3!}{1 + 2 + 2^2/2! + 2^3/3!}
  = \frac{1.333}{6.333} = 0.211 .
\end{equation*}
رابطهٔ بازگشتی همین را می‌دهد: $B_1 = 2/3$، $B_2 = (2\cdot\tfrac23)/(2+\tfrac43) = 0.4$ و
$B_3 = 0.8/(3+0.8) = 0.211$.

\textbf{تفسیر.} از هر پنج تماس یکی از دست می‌رود: یک سایت سه‌کاناله در \lr{2\,E} بی‌فایده است.
مخزن‌های کوچک برای درجهٔ خدمت قابل قبول به بار بسیار کم نیاز دارند.
\end{examplebox}

\begin{examplebox}{نرخ ورود در ساعت اوج}\label{ex:lambda}
\textbf{مسئله.} میانگین مدت تماس ۹۰ ثانیه است. چه نرخ ورودی \lr{82.2\,E} \cref{ex:dim} را تولید
می‌کند و \acro{MSC} در آن سلول باید چند برقراری تماس در ثانیه را مدیریت کند؟

\textbf{حل.} $\lambda = A/h = 82.2/90 = 0.913$ تماس بر ثانیه، یعنی حدود ۳۲۹۰ تلاش تماس در ساعت
در یک سلول. هر تلاش ده‌ها پیام سیگنالینگ (\lr{RR, MM, CC, BSSMAP, ISUP}) تولید می‌کند؛ پس بار
سیگنالینگ با $\lambda$ مقیاس می‌شود، نه با $A$.

\textbf{تفسیر.} تماس‌های کوتاه (یا پیامک و به‌روزرسانی مکان که از \acro{SDCCH} استفاده می‌کنند)
ظرفیت سیگنالینگ را تحت فشار قرار می‌دهند که باید جداگانه ابعادبندی شود (\cref{exr:sdcch}).
\end{examplebox}

% ---------------------------------------------------------------------
\section{مقایسهٔ فناوری‌ها}
\begin{table}[htbp]
\centering\small
\caption{سوئیچ مداری در برابر سوئیچ بسته‌ای.}\label{tab:csps}
\rowcolors{2}{acnmist}{white}
\begin{tabularx}{\linewidth}{@{}l Y Y@{}}
\toprule
\rowcolor{white} & سوئیچ مداری (صدای \acro{GSM}) & سوئیچ بسته‌ای (\acro{GPRS} و پس از آن) \\
\midrule
تخصیص منابع & هنگام برقراری و برای کل تماس رزرو می‌شود & به ازای هر بسته یا هر بازهٔ زمان‌بندی \\
تأخیر       & ثابت و کم                                  & متغیر (صف) \\
سازوکار اتلاف & انسداد هنگام برقراری                     & اتلاف یا تأخیر بسته تحت بار \\
بازدهی      & برای منابع انفجاری کم                       & زیاد (تسهیم آماری) \\
مدل ابعادبندی & \lr{Erlang-B} ($M/M/m/m$)                & مدل‌های صف ($M/M/1$ و غیره) \\
صورت‌حساب   & زمان                                       & حجم یا نرخ ثابت \\
سیگنالینگ   & \acro{SS7}: \lr{ISUP + MAP}                & \acro{GTP} و سپس \lr{Diameter} و \lr{HTTP/2} \\
\bottomrule
\end{tabularx}
\end{table}

\paragraph{چرا سوئیچ مداری برای صدا درست بود} صدا جریانی با نرخ ثابت و حساس به تأخیر است؛ رزرو
\qty{13}{\kilo\bit\per\second} کیفیت را بدون زمان‌بندی پیچیده تضمین می‌کند و مدل ارلانگ امکان
ابعادبندی دقیق را می‌دهد. \paragraph{چرا به محدودیت تبدیل شد} داده انفجاری است: یک صفحهٔ وب
یک ثانیه نرخ بالا و یک دقیقه هیچ نیاز دارد. رزرو مدار برای آن اتلاف است و یک مدار
\qty{9.6}{\kilo\bit\per\second} بسیار کند است. اجرای دو شبکهٔ هسته (مداری و بسته‌ای) هزینهٔ
بهره‌برداری را نیز دو برابر می‌کرد -- انگیزهٔ هستهٔ تمام‌\acro{IP} در \acro{LTE} (درس ۵)، جایی که
خود صدا نیز به بسته تبدیل می‌شود (\acro{VoLTE}).

% ---------------------------------------------------------------------
\section{نکته‌های کلیدی}
\begin{takeawaybox}
\begin{itemize}
  \item یک تماس \acro{GSM} یک شیار رادیویی، کانال‌های \lr{Abis/A}، یک مسیر سوئیچ و ترانک‌ها را
        برای کل مدتش رزرو می‌کند -- حتی در سکوت.
  \item \acro{SS7} سیگنالینگ را از صدا جدا می‌کند: \acro{MTP} انتقال می‌دهد، \acro{SCCP} نشانی‌دهی
        می‌کند (عنوان‌های سراسری)، \acro{TCAP} گفت‌وگوها را ساختار می‌دهد، \acro{MAP} تحرک را
        مدیریت می‌کند و \acro{ISUP} ترانک‌ها را کنترل می‌کند.
  \item یک تماس ورودی پیش از مسیریابی و فراخوانی به \acro{HLR} (\msg{SRI}) و \acro{VLR} میزبان
        (\msg{PRN}) برای گرفتن یک \acro{MSRN} نیاز دارد.
  \item \lr{Erlang-B}، یعنی $B(A,m)$، سامانه‌های اتلاف را ابعادبندی می‌کند؛ آن را با رابطهٔ بازگشتی
        \cref{eq:erlangrec} محاسبه کنید.
  \item مخزن‌های بزرگ‌تر کاراترند (بهرهٔ ترانکینگ)؛ تقسیم مخزن‌ها ظرفیت را کم می‌کند.
  \item سوئیچ مداری بخش ساکت هر تماس را هدر می‌دهد و نمی‌تواند به‌طور کارا به دادهٔ انفجاری
        خدمت دهد؛ همین، حرکت به سوئیچ بسته‌ای را رقم زد.
\end{itemize}
\end{takeawaybox}

% ---------------------------------------------------------------------
\section{تمرین‌ها}

\begin{exercisebox}{1}{سیگنالینگ}\label{exr:ss7}
چرا \acro{ISUP} مستقیماً روی \lr{MTP 3} اجرا می‌شود ولی \acro{MAP} به \acro{SCCP} نیاز دارد؟
پاسخ خود را به نحوهٔ معلوم بودن مقصد هر پیام ربط دهید.
\end{exercisebox}

\begin{exercisebox}{1}{شماره‌ها}\label{exr:msrn}
چرا \acro{GMSC} نمی‌تواند تماس ورودی را با \acro{IMSI} مسیریابی کند؟ چرا \acro{MSRN} به محض
رسیدن تماس به \acro{MSC} خدمت‌دهنده آزاد می‌شود؟
\end{exercisebox}

\begin{exercisebox}{2}{ابعادبندی}\label{exr:dim}
یک سلول ۳۰ کانال ترافیک دارد. (الف) در انسداد ۱٪ و ۵٪ چه ترافیکی را می‌تواند بپذیرد؟
(ب) با \lr{20\,mE} برای هر مشترک، در ۵٪ چند مشترک بیشتر از ۱٪ خدمت می‌گیرند؟
\end{exercisebox}

\begin{exercisebox}{2}{ظرفیت \acro{SDCCH}}\label{exr:sdcch}
به‌روزرسانی‌های مکان، پیامک‌ها و برقراری تماس‌ها هر یک یک \acro{SDCCH} را حدود ۳ ثانیه اشغال
می‌کنند. یک سلول در ساعت اوج ۱٫۵ تراکنش از این نوع در ثانیه می‌بیند. برای انسداد حداکثر
۰٫۵٪ چند زیرکانال \acro{SDCCH} لازم است؟
\end{exercisebox}

\begin{exercisebox}{3}{تلاش مجدد}\label{exr:retry}
فرض کنید هر تماس‌گیرندهٔ مسدودشده یک بار، بی‌درنگ و مستقل دوباره تلاش می‌کند. نشان دهید
ترافیک ارائه‌شده تقریباً $A' = A(1+P_{\mathrm{B}}')$ می‌شود که در آن
$P_{\mathrm{B}}' = B(A',m)$؛ این معادلهٔ نقطهٔ ثابت را برای $m = 30$ و $A = 21.93$\,E به‌صورت
عددی حل کنید و نتیجه را تفسیر کنید.
\end{exercisebox}

\begin{exercisebox}{3}{مسیر رفت‌وبرگشتی}\label{exr:trombone}
مشترک یک اپراتور ایرانی در ترکیه رومینگ است و یک کاربر تلفن ثابت ترکیه با او تماس می‌گیرد.
مسیر مدار را رسم کنید و بنویسید هر بخش آن را کدام شبکه حمل می‌کند. یک تغییر سیگنالینگی
پیشنهاد کنید که از بخش‌های بین‌المللی جلوگیری کند و مشخص کنید شبکهٔ ترکیه به چه اطلاعاتی
نیاز دارد.
\end{exercisebox}

\subsection*{پاسخ‌ها}
\addcontentsline{toc}{subsection}{پاسخ‌ها}

\begin{solutionbox}{\ref{exr:ss7}}
پیام‌های \acro{ISUP} میان مراکز مجاور و در امتداد همان مداری که در حال برقراری است حرکت می‌کنند؛
مرکز بعدی و کد نقطهٔ آن از جدول‌های مسیریابی معلوم است، پس مسیریابی \lr{MTP 3} کافی است. پیام‌های
\acro{MAP} به پایگاه‌های داده‌ای نشانی‌دهی می‌شوند که با هویت مشترک شناخته می‌شوند (مثلاً «\acro{HLR}
مربوط به \lr{IMSI 432\ldots}») و ممکن است در کشور دیگری باشند؛ ترجمهٔ عنوان سراسری \acro{SCCP}
چنین هویتی را گام‌به‌گام به یک کد نقطه تبدیل می‌کند.
\end{solutionbox}

\begin{solutionbox}{\ref{exr:msrn}}
\acro{PSTN} با شماره‌های \lr{E.164} مسیریابی می‌کند؛ \acro{IMSI} شمارهٔ قابل شماره‌گیری نیست و
عمداً بیرون از شبکهٔ موبایل افشا نمی‌شود. \acro{MSRN} شمارهٔ موقت و کمیابی از مخزن \acro{MSC}
میزبان است؛ وقتی تماس به آن \acro{MSC} رسید، \acro{IMSI} از \acro{VLR} بازیابی می‌شود و \acro{MSRN}
برای تماس دیگری قابل استفادهٔ دوباره است.
\end{solutionbox}

\begin{solutionbox}{\ref{exr:dim}}
(الف) $B(A,30)=0.01 \Rightarrow A = 20.3$\,E و $B(A,30)=0.05 \Rightarrow A = 24.8$\,E.
(ب) مشترکان: $20.3/0.02 \approx 1017$ در برابر $24.8/0.02 \approx 1240$، یعنی حدود ۲۲۰ مشترک
بیشتر (۲۲٪) به بهای پنج برابر شدن انسداد.
\end{solutionbox}

\begin{solutionbox}{\ref{exr:sdcch}}
$A = \lambda h = 1.5\times3 = 4.5$\,E. با تکرار \cref{eq:erlangrec}: $B(4.5,10)=0.0105$ و
$B(4.5,11)=0.0043$. یازده زیرکانال لازم است؛ در عمل اپراتورها دو کانال \lr{SDCCH/8} (۱۶ زیرکانال)
پیکربندی می‌کنند که جا برای اوج‌ها نیز باقی می‌گذارد.
\end{solutionbox}

\begin{solutionbox}{\ref{exr:retry}}
کسری $P_{\mathrm{B}}'$ از $\lambda$ تلاش نخست یک تلاش دوم تولید می‌کند؛ پس نرخ کل تلاش‌ها
$\lambda(1+P_{\mathrm{B}}')$ و $A' = A(1+P_{\mathrm{B}}')$ است (با فرض پواسون بودن تلاش‌های مجدد که
یک تقریب است). با تکرار از $P_{\mathrm{B}}'=0.02$:
$A' = 22.37$، $B = 0.0236$؛ $A' = 22.45$، $B = 0.0243$؛ و همگرایی به
$A' \approx 22.47$\,E و $P_{\mathrm{B}}' \approx 0.0244$. انسدادی که هر تلاش می‌بیند حدود ۲۲٪
افزایش می‌یابد؛ کسری از \emph{کاربران} که در نهایت ناموفق می‌مانند
$P_{\mathrm{B}}'^{\,2}\approx 0.06\%$ است، اما شبکه سیگنالینگ بیشتری حمل می‌کند. تلاش‌های مجدد یک
بازخورد مثبت‌اند که می‌تواند شبکهٔ پربار را ناپایدار کند.
\end{solutionbox}

\begin{solutionbox}{\ref{exr:trombone}}
مسیر: تماس‌گیرندهٔ ترکیه‌ای $\to$ \acro{PSTN} ترکیه $\to$ بخش بین‌المللی تا \acro{GMSC} خانگی در
ایران (\acro{MSISDN} ایرانی است) $\to$ پرسش از \acro{HLR} یک \acro{MSRN} ترکیه‌ای برمی‌گرداند $\to$
بخش بین‌المللی دوم به \acro{MSC} میزبان در ترکیه $\to$ مشترک رومینگ. مسیریابی بهینه به دروازهٔ
ترکیه اجازه می‌دهد مستقیماً از \acro{HLR} خانگی بپرسد (\acro{MAP} \acro{SRI} به سوی \acro{HPLMN})
و تماس را درون ترکیه به \acro{MSRN} مسیریابی کند؛ این کار مستلزم آن است که شبکهٔ میزبان شماره را
متعلق به یک مشترک رومینگ تشخیص دهد و توافقی برای پرسش از \acro{HLR} وجود داشته باشد، و صورت‌حساب
را تغییر می‌دهد؛ به همین دلیل به‌ندرت پیاده شد.
\end{solutionbox}

% ---------------------------------------------------------------------
\section{مطالعهٔ بیشتر}
\acro{SS7} در \cite{ituq700} معرفی شده است؛ \acro{MAP} در \cite{ts29002}، مدیریت مکان در
\cite{ts23012} و مدیریت تماس در \cite{ts23018} مشخص شده‌اند. پیام‌های کنترل تماس و تحرک لایهٔ ۳
در \cite{ts24008} آمده‌اند. \lr{Kleinrock} \cite{kleinrock1975queueing} فرمول اتلاف را به‌طور دقیق
استخراج می‌کند و \lr{Rappaport} \cite{rappaport2002wireless} آن را در ابعادبندی سلولی به کار
می‌برد. \lr{Eberspächer} و همکاران \cite{eberspaecher2009gsm} جریان‌های تماس \acro{GSM} را با
جزئیات توصیف می‌کنند.

\begin{latin}
\printbibliography[heading=subbibliography,title={References}]
\end{latin}
\end{document}
